\subsection{Valuations of height 1}

PROPOSITION 6. Let K be a field and A a subring of K. Suppose that A is not a field. Then the following conditions are equivalent:

\begin{enumerate}
\def\labelenumi{(\alph{enumi})}
\item
  A is the ring of a valuation of height 1 on K;
\item
  A is a valuation ring of K and has no prime ideals other than (0) and m(A);
\item
  A is maximal among the subrings of K distinct from K.
\end{enumerate}

Proposition 5 of no. 4 shows that (a) implies (b) and Proposition 1 of no. 1 shows that (b) implies (c). It remains to show that (c) implies (a). Suppose A is maximal among the subrings of K distinct from K. Let m be a maximal ideal of A and V a valuation ring of K dominating \(A_{m}\) (§ 1, no. 2, Corollary to Theorem 2); as \(m(V) \cap A = m\) and \(m \neq (0)\) (since A is not a field), \(V \neq K\) , whence V = A, which proves that A is the ring of a valuation v on K. This being so, v is of height 1 by Propositions 1 (no. 1) and 5 (no. 4).

PROPOSITION 7. For a valuation on a field to be of height 1, it is necessary and sufficient that its order group be isomorphic to a non-zero ordered subgroup of R.

This follows in fact from the following proposition:

PROPOSITION 8. Let G be a totally ordered group not reduced to 0. The following conditions are equivalent:

\begin{enumerate}
\def\labelenumi{(\alph{enumi})}
\item
  G is of height 1;
\item
  for all \(x > 0\) and \(y \geqslant 0\) in \(G\) , there exists an integer \(n \geqslant 0\) such that \(y \leqslant nx\) ;
\item
  G is isomorphic to a subgroup of the ordered additive group R which is not reduced to 0.
\end{enumerate}

Let \(x\) be a positive element of \(G\) and let \(H_x\) be the set of \(y \in G\) such that there exists an integer \(n \geqslant 0\) satisfying \(|y| \leqslant nx\) . It is easily verified that \(H_x\) is an isolated subgroup of \(G\) and that every isolated subgroup of \(G\) containing \(x\) contains \(H_x\) . Condition (a) is therefore equivalent to `` \(H_x = G\) for all \(x > 0\) '', that is, to condition (b).

Clearly (c) implies (b). Conversely, suppose condition (b) holds and let Q denote the set of elements >0 of G. Suppose first that Q has a least element x; for all \(y \in Q\) , let n be the least integer such that \(y \leqslant nx\) ; if y < nx, then also \(nx - y \geqslant x\) , whence \(y \leqslant (n - 1)x\) contrary to the choice of n; then y = nx, which shows that \(G = Zx\) is isomorphic to \(Z \subset R\) . Suppose now that Q has no least element; we apply to the ordered set \(P = Q \cup \{0\}\) Proposition 1 of General Topology, Chapter V, §2 (which is possible, since condition (b) is just ``Archimedes' axiom''); it is seen that there exists a strictly increasing mapping f of P to \(R_{+}\) such that

\[
f (x + y) = f (x) + f (y)
\]

for \(x \in P\) and \(y \in P\) ; by linearity f can be extended to an isomorphism of G onto a subgroup of R, which proves that (b) implies (c).

PROPOSITION 9. Let K be a field, v a non-improper valuation on K and A the ring of v. For A to be completely integrally closed (Chapter V, § 1, no. 4, Definition 5), it is necessary and sufficient that v be of height 1.

Suppose \(v\) is of height 1. Let \(x \in \mathbf{K}\) be such that the \(x^n (n \geqslant 0)\) are all contained in a finitely generated sub-A-module of \(\mathbf{K}\) . There exists \(d \in \mathbf{A} - \{0\}\) such that \(dx^n \in \mathbf{A}\) for all \(n \geqslant 0\) . Then \(v(d) + nv(x) \geqslant 0\) , that is \(n(-v(x)) \leqslant v(d)\) for all \(n \geqslant 0\) , whence \(-v(x) \leqslant 0\) (Proposition 8 (b)) and \(x \in \mathbf{A}\) . Thus \(\mathbf{A}\) is completely integrally closed.

Suppose now that \(v\) is not of height 1. Then there exist \(y \in \mathfrak{m}(\mathbf{A})\) and \(t \in \mathbf{A}\) such that \(nv(y) < v(t)\) for all \(n \geqslant 0\) (Proposition 8 (b)). Then \(ty^{-n} \in \mathbf{A}\) for all \(n \geqslant 0\) , but \(y^{-1} \notin \mathbf{A}\) . Hence \(\mathbf{A}\) is not completely integrally closed.

COROLLARY. Let \(\mathbf{K}\) be a field, \((v_{\alpha})_{\alpha \in \mathbf{I}}\) a family of valuations of height 1 on \(\mathbf{K}\) and \(\mathbf{A}\) the intersection of the rings of the \(v_{\alpha}\) . Then \(\mathbf{A}\) is a completely integrally closed domain.

A completely integrally closed domain is not always an intersection of valuation rings of height 1 (Exercise 6).

